% SOURCE_AUTHORITY: sga5_fr_workpass.tex, lines 8903--8946 (LF-normalized unit).
% SOURCE_SHA256: 791F4EFFC5E02832D5D77ED03518C8156D6F07E4C8238B03545DB93D883FBB28
\subsubsection*{2.2.3. Outras propriedades}

\paragraph{A) Dimensão cohomológica.}
Seja \(f:X\to Y\) um morfismo separado de tipo finito de pré-esquemas localmente noetherianos. Supõe-se que a dimensão relativa de \(f\) é menor ou igual a um inteiro \(d\). Então, para todo \((AR,\mathbb Z_\ell)\)-feixe construtível \(F\),
\[
\R^i_!f(F)=0\qquad\text{para }i>2d.
\]
Isso se segue da afirmação análoga para as componentes de \(F\) (SGA 4 X 4.1 e 5.2).

\paragraph{B) Mudança de base própria.}
Consideremos um quadrado cartesiano
\[
\begin{tikzcd}
X' \arrow[r,"g'"] \arrow[d,"f'"'] & X \arrow[d,"f"]\\
Y' \arrow[r,"g"] & Y .
\end{tikzcd}
\]
Sejam \(f:X\to Y\) e \(g:Y'\to Y\) morfismos de pré-esquemas localmente noetherianos, com \(f\) separado do tipo finito. Seja \(X'=X\times_Y Y'\), que é localmente noetheriano já que \(f\) é do tipo finito, e sejam
\[
f':X'\to Y',\qquad g':X'\to X
\]
os morfismos canônicos. Para todo inteiro \(i\) e todo \(\mathbb Z_\ell\)-feixe construtível \(F\), existe um isomorfismo funtorial em \(F\) em \(\mathbb Z_\ell\text{-}\fc(Y')\):
\[
w_i:g^*(\R^i_!f(F))\longrightarrow \R^i_!f'(g'^*F).
\]
Além disso, se
\[
0\longrightarrow F'\longrightarrow F\longrightarrow F''\longrightarrow 0
\]
é uma sucessão exata de \(\mathbb Z_\ell\text{-}\fc(X)\) (resp. \((AR,\mathbb Z_\ell)\text{-}\fc(X)\)), os \((w_i)_{i\in\mathbb N}\) comutam com os operadores de bordo; em outras palavras, os diagramas que aparecem a seguir são comutativos:
\[
\begin{tikzcd}[column sep=huge]
g^*(\R^i_!f(F'')) \arrow[r,"w_i"] \arrow[d,"g^*(\delta)"']
&
\R^i_!f'(g'^*F'') \arrow[d,"\delta"]\\
g^*(\R^{i+1}_!f(F')) \arrow[r,"w_{i+1}"]
&
\R^{i+1}_!f'(g'^*F').
\end{tikzcd}
\]

\begin{prooftext}[Demonstração]
Basta demonstrar a afirmação análoga para os \((AR,\mathbb Z_\ell)\)-feixes construtíveis. Ora, ela se deduz, por passagem ao quociente, da afirmação análoga para os sistemas projetivos de feixes de torção, que é consequência imediata do teorema de mudança de base própria (cf. SGA 4 XII 5.1 no caso próprio e SGA 4 XVII no caso geral).
\end{prooftext}
